\section*{Chapter \ref{ch:b3:rings} --- Rings and Arithmetic}

\begin{solution}{exo:b3:rings:1}
(a) Evaluation $f \colon \Z[X] \to \Z[\iu]$, $P \mapsto P(\iu)$,
is a surjective ring morphism ($a + bX \mapsto a + b\iu$). Kernel:
divide $P$ by the \emph{monic} $X^2 + 1$ in $\Z[X]$: $P = (X^2 +
1)Q + (bX + a)$ with $a, b \in \Z$; then $P(\iu) = a + b\iu = 0$
iff $a = b = 0$. So $\ker f = (X^2+1)$ and
\cref{thm:b3:rings:firstiso} concludes.

(b) The same computation with $\R$-coefficients: $\R[X]/(X^2+1)
\cong \C$ --- this is the cleanest \emph{construction} of $\C$.

(c) $X^2 + X + 1$ has no root in $\mathbb F_2$ ($0, 1 \mapsto 1$),
so, having degree $2$, it is irreducible: the quotient $\mathbb
F_4 = \mathbb F_2[X]/(X^2+X+1)$ is a field
(\cref{prop:b3:rings:primemaximal}; $(P)$ maximal in $K[X]$ when
$P$ is irreducible, since $K[X]$ is a PID: an ideal $(D) \supseteq
(P)$ means $D \mid P$). Its four elements are $0, 1, \omega,
\omega + 1$ where $\omega = \bar X$, with $\omega^2 = \omega + 1$.
Multiplication table (nonzero elements):
\[
\omega \cdot \omega = \omega + 1, \qquad
\omega(\omega + 1) = \omega^2 + \omega = 1, \qquad
(\omega+1)^2 = \omega^2 + 1 = \omega .
\]
The nonzero elements form a cyclic group of order $3$ generated by
$\omega$.
\end{solution}

\begin{solution}{exo:b3:rings:2}
(a) Let $p$ be irreducible in a UFD and $p \mid ab$, say $ab =
pc$, with $a, b \neq 0$ (else trivial). If $a$ or $b$ is a unit,
$p$ divides the other. Otherwise factor $a$, $b$ and $c$ into
irreducibles: the two factorizations of $ab$,
\[
(\text{factors of } a)(\text{factors of } b) = p \cdot
(\text{factors of } c),
\]
must agree up to order and associates: $p$ is an associate of some
irreducible factor of $a$ or of $b$, hence divides it.

(b) Let $A$ be a finite domain and $x \neq 0$. The map $y \mapsto
xy$ is injective ($xy = xy' \Rightarrow x(y - y') = 0 \Rightarrow
y = y'$), hence surjective ($A$ finite): $1 = xy$ for some $y$.

(c) If $\mathfrak p$ is prime in a finite ring $A$, then
$A/\mathfrak p$ is a finite domain, hence a field by (b), so
$\mathfrak p$ is maximal (\cref{prop:b3:rings:primemaximal}).
\end{solution}

\begin{solution}{exo:b3:rings:3}
Norms: $N(3) = 9$, $N(1 \pm \iu\sqrt5) = 6$, $N(2 \pm \iu\sqrt5) =
9$. The equations $x^2 + 5y^2 = 2$ and $x^2 + 5y^2 = 3$ have no
integer solutions, so no element has norm $2$ or $3$. A proper
factorization of $3$ would need two factors of norm $3$:
impossible --- $3$ is irreducible. A proper factorization of $1
\pm \iu\sqrt5$ (norm $6$) would need factors of norms $2, 3$:
impossible. Same for $2 \pm \iu\sqrt5$ (norm $9$: factors would
have norm $3$). Now
\[
9 = 3 \cdot 3 = (2 + \iu\sqrt5)(2 - \iu\sqrt5),
\]
two factorizations into irreducibles. They are genuinely
different: the units are $\pm 1$ (norm $1$), and $2 \pm \iu\sqrt5
\neq \pm 3$. So uniqueness fails --- while existence of
factorizations holds in $\Z[\iu\sqrt5]$
(\cref{exo:b3:rings:10}(c)): non-factoriality here is purely a
uniqueness failure. (Consistently, \cref{prop:b3:rings:primeirred}:
these irreducibles are not prime.)
\end{solution}

\begin{solution}{exo:b3:rings:4}
(a) Let $a, b \in \Z[\iu]$, $b \neq 0$, and $a/b = x + \iu y \in
\C$. Choose integers $m, n$ with $\abs{x - m} \leq \frac12$,
$\abs{y - n} \leq \frac12$, and set $q = m + \iu n$, $r = a - bq$.
Then
\[
N(r) = N(b)\,\abs*{\tfrac ab - q}^2
\leq N(b)\Bigl(\tfrac14 + \tfrac14\Bigr) = \tfrac{N(b)}2 < N(b).
\]
So $N$ is a Euclidean function ($N(r) < N(b)$ or $r = 0$).

(b) If $uv = 1$ then $N(u)N(v) = 1$ with $N(u) \in \N$: $N(u) =
1$, i.e.\ $x^2 + y^2 = 1$: $u \in \{\pm 1, \pm\iu\}$; conversely
these are units.

(c) For $\Z[\iu\sqrt2]$: the same rounding gives $\abs{a/b - q}^2
\leq \frac14 + \frac{2}4 = \frac34 < 1$: Euclidean; units: $x^2 +
2y^2 = 1$ gives $\pm 1$. For $\Z[\iu\sqrt5]$ the bound becomes
$\frac14 + \frac54 = \frac32 > 1$: the rounding argument fails ---
and must fail, since $\Z[\iu\sqrt5]$ is not even a UFD
(\cref{exo:b3:rings:3}), while Euclidean would imply UFD
(\cref{thm:b3:rings:euclideanpid,thm:b3:rings:pidufd}).
\end{solution}

\begin{solution}{exo:b3:rings:5}
(a) If $x^n = 0$:
\[
(1 + x)\bigl(1 - x + x^2 - \dots + (-1)^{n-1}x^{n-1}\bigr)
= 1 + (-1)^{n-1}x^n = 1 .
\]
(b) In a domain, $\deg(PQ) = \deg P + \deg Q$; $PQ = 1$ forces
$\deg P = \deg Q = 0$ and $P, Q \in A^\times$: $A[X]^\times =
A^\times$. Over $\Z/4\Z$: $(1 + 2X)^2 = 1 + 4X + 4X^2 = 1$, so $1
+ 2X$ is a unit of degree $1$ (here $2$ is nilpotent; compare
(a)).

(c) $e^2 = e$ gives $e(e - 1) = 0$, so $e \in \{0, 1\}$ in a
domain. If $x^n = 0$ with $n \geq 1$ minimal and $x \ne 0$, then
$n \geq 2$ and $x \cdot x^{n-1} = 0$ with both factors nonzero:
contradiction.
\end{solution}

\begin{solution}{exo:b3:rings:6}
(a) $K[X,Y]/(X,Y) \cong K$ (evaluate at $(0,0)$): a field, so
$(X,Y)$ is maximal. If $(X, Y) = (P)$: $P \mid X$ forces (degrees
in $Y$) $P \in K[X]$, and $P \mid Y$ then forces $P \in K$; $P =
0$ is absurd and $P \in K^\times$ would give $(P) = K[X,Y]$,
contradicting properness ($K[X,Y]/(X,Y) \cong K \neq 0$). So
$(X,Y)$ is not principal.

(b) Evaluation $P(X, Y) \mapsto P(X, X^2)$ maps $K[X,Y]$ onto
$K[X]$; its kernel is $(Y - X^2)$: dividing by the monic-in-$Y$
polynomial $Y - X^2$, $P = (Y - X^2)Q + R(X)$, and $P(X, X^2) =
R(X)$. So $K[X,Y]/(Y - X^2) \cong K[X]$ --- the coordinate ring of
a parabola, isomorphic to a line's.

Evaluation $P(X, Y) \mapsto P(X, X^{-1})$ maps $K[X, Y]$ onto the
ring $K[X, X^{-1}]$ of Laurent polynomials. Its kernel contains
$(XY - 1)$; conversely, modulo $XY - 1$ every class has a
representative $R = \sum_{n \geq 0} a_nX^n + \sum_{m \geq 1} b_m
Y^m$ (replace each product $XY$ by $1$ repeatedly), and $R(X,
X^{-1}) = \sum a_n X^n + \sum b_m X^{-m} = 0$ forces all $a_n =
b_m = 0$. Hence $K[X, Y]/(XY - 1) \cong K[X, X^{-1}]$ --- the
coordinate ring of a hyperbola: the line with one point removed.

(c) $(Y - X^2)$ is prime (the quotient $K[X]$ is a domain) but not
maximal ($K[X]$ is not a field; concretely $(Y - X^2) \subsetneq
(Y - X^2,\, X) \subsetneq K[X,Y]$).
\end{solution}

\begin{solution}{exo:b3:rings:7}
\emph{$X^5 - 12X^3 + 36X - 12$}: Eisenstein at $p = 3$ ($3 \mid
12, 36, 12$; $9 \nmid 12$; $3 \nmid 1$): irreducible. (At $p = 2$
Eisenstein fails: $4 \mid 12$.)

\emph{$X^4 + X + 1$}: reduce mod $2$. No root in $\mathbb F_2$;
the only irreducible quadratic over $\mathbb F_2$ is $X^2 + X +
1$, and $(X^2+X+1)^2 = X^4 + X^2 + 1 \neq X^4 + X + 1$. So $X^4 +
X + 1$ is irreducible over $\mathbb F_2$, hence over $\Q$
(\cref{thm:b3:rings:criteria}(1); it is monic).

\emph{$X^4 + 4$}: reducible --- the Sophie Germain identity,
$X^4 + 4 = (X^2 - 2X + 2)(X^2 + 2X + 2)$.

\emph{$X^4 + 1$}: shift, $(X+1)^4 + 1 = X^4 + 4X^3 + 6X^2 + 4X +
2$: Eisenstein at $2$. A factorization of $X^4+1$ would shift to
one of $(X+1)^4 + 1$: irreducible.

\emph{$X^3 - X - 1$}: a cubic is reducible over $\Q$ iff it has a
rational root; a rational root of a monic integer polynomial is an
integer dividing the constant term (rational root theorem: if
$(p/q)$ in lowest terms is a root, $q \mid 1$, $p \mid -1$), and
$\pm 1$ are not roots ($-1$ and $-1$): irreducible.
\end{solution}

\begin{solution}{exo:b3:rings:8}
(a) For $\gcd(m, n) = 1$, \cref{thm:b3:rings:crt} gives a ring
isomorphism $\Z/mn\Z \cong \Z/m\Z \times \Z/n\Z$. An element of a
product ring is a unit iff both coordinates are, so
$(\Z/mn\Z)^\times \cong (\Z/m\Z)^\times \times (\Z/n\Z)^\times$
and $\varphi(mn) = \varphi(m)\varphi(n)$. For a prime power, the
non-units of $\Z/p^k\Z$ are the classes of multiples of $p$:
$\varphi(p^k) = p^k - p^{k-1}$. Hence
\[
\varphi(n) = \prod_i \bigl(p_i^{\alpha_i} -
p_i^{\alpha_i-1}\bigr) = n \prod_{p \mid n}\Bigl(1 -
\frac1p\Bigr).
\]

(b) $M = 7 \cdot 11 \cdot 13 = 1001$. Idempotents: $e_1 \equiv
(1, 0, 0)$: $143 = 11\cdot13 \equiv 3 \pmod 7$ and $3 \cdot 5
\equiv 1$: $e_1 = 143 \cdot 5 = 715$. $e_2$: $91 \equiv 3
\pmod{11}$, $3 \cdot 4 \equiv 1$: $e_2 = 91\cdot4 = 364$. $e_3$:
$77 \equiv -1 \pmod{13}$: $e_3 = 77 \cdot 12 = 924$. Then
\[
x \equiv 2\,e_1 + 5\,e_2 + 1\,e_3 = 1430 + 1820 + 924 = 4174
\equiv 170 \pmod{1001},
\]
and indeed $170 = 24\cdot7 + 2 = 15\cdot11 + 5 = 13\cdot13 + 1$.
\end{solution}

\begin{solution}{exo:b3:rings:9}
(a) If $x^n = 0$ and $y^m = 0$, the binomial expansion of
$(x+y)^{n+m}$ has every term $x^iy^j$ with $i + j = n + m$, so $i
\geq n$ or $j \geq m$: each term vanishes, and $x + y$ is
nilpotent; $(ax)^n = a^nx^n = 0$: $\operatorname{Nil}(A)$ is an
ideal. If $\mathfrak p$ is prime and $x^n = 0 \in \mathfrak p$,
induction on $n$ gives $x \in \mathfrak p$ ($x \cdot x^{n-1} \in
\mathfrak p$).

(b) Let $a \notin \operatorname{Nil}(A)$ and $S = \{a^n : n \geq
1\}$, so $0 \notin S$. The set $\mathcal E$ of ideals disjoint
from $S$ contains $(0)$ and is inductive (the union of a chain of
ideals disjoint from $S$ is an ideal disjoint from $S$): Zorn
provides $\mathfrak p \in \mathcal E$ maximal. $\mathfrak p$ is
proper ($a \notin \mathfrak p$, since $a \in S$). Primality: let
$x, y \notin \mathfrak p$. By maximality, $\mathfrak p + (x)$ and
$\mathfrak p + (y)$ meet $S$: $a^m \in \mathfrak p + (x)$, $a^n
\in \mathfrak p + (y)$. Multiplying, $a^{m+n} \in \mathfrak p +
(xy)$. If $xy \in \mathfrak p$, then $a^{m+n} \in \mathfrak p \cap
S$: absurd. So $xy \notin \mathfrak p$ --- contrapositive of
primality. Hence every non-nilpotent element avoids some prime
ideal; with (a),
$\operatorname{Nil}(A) = \bigcap_{\mathfrak p} \mathfrak p$.
\end{solution}

\begin{solution}{exo:b3:rings:10}
(a) The chain $\ker f \subseteq \ker f^2 \subseteq \cdots$
stabilizes (\cref{prop:b3:rings:noethacc}): $\ker f^n = \ker
f^{n+1}$ for some $n$. Let $x \in \ker f$. As $f$, hence $f^n$, is
surjective, $x = f^n(y)$ for some $y$; then $f^{n+1}(y) = f(x) =
0$, so $y \in \ker f^{n+1} = \ker f^n$, i.e.\ $x = f^n(y) = 0$.

(b) $I_n = \{f \in \mathcal C(\intcc01, \R) : f\restriction_{
\intcc0{1/n}} = 0\}$ is an ideal, and $I_n \subseteq I_{n+1}$. The
inclusion is strict: $x \mapsto \max\bigl(0, x -
\frac1{n+1}\bigr)$ vanishes on $\intcc0{\frac1{n+1}}$ but not on
$\intcc0{\frac1n}$. An infinite strictly increasing chain
contradicts \cref{prop:b3:rings:noethacc}.

(c) Suppose the set of nonzero nonunits admitting no factorization
into irreducibles is nonempty. The corresponding family of ideals
$\{(a)\}$ has a maximal element $(a)$
(\cref{prop:b3:rings:noethacc}). The element $a$ is not
irreducible (an irreducible is its own factorization), so $a = bc$
with $b, c$ nonunits; $(a) \subseteq (b)$ is strict (as $(a) =
(b)$ would give $b = ad$, $a = adc$, so $dc = 1$: $c$ a unit),
likewise $(a) \subsetneq (c)$. By maximality, $b$ and $c$ both
factor into irreducibles; concatenating factors $a$:
contradiction. Applied to $\Z[\iu\sqrt5]$ --- Noetherian as a
quotient of $\Z[X]$ (\cref{thm:b3:rings:hilbert}, $\Z[\iu\sqrt5]
\cong \Z[X]/(X^2+5)$) --- this shows factorizations exist there;
\cref{exo:b3:rings:3} showed uniqueness is what fails.
\end{solution}

\begin{solution}{exo:b3:rings:11}
(a) \cref{exo:b3:rings:7}: shift and Eisenstein at $2$.

(b) Mod $2$: $X^4 + 1 = (X + 1)^4$. Now let $p$ be odd. The
squares form a subgroup of index $2$ in $(\Z/p\Z)^\times$: the
morphism $x \mapsto x^2$ has kernel $\{\pm 1\}$ (two elements: $X^2
- 1$ has at most $2$ roots in a field, and $1 \neq -1$ for odd
$p$), so its image has $\frac{p-1}2$ elements. Consequently, the
product of two non-squares is a square (in the order-$2$ quotient
group, $\overline{xy} = \bar x\bar y$). Hence \emph{at least one
of $-1$, $2$, $-2$ is a square mod $p$} (if $-1$ and $2$ are not,
$-2 = (-1)\cdot 2$ is). In each case $X^4 + 1$ factors mod $p$:
\begin{itemize}
  \item $-1 = c^2$: \quad $X^4 + 1 = X^4 - c^2 = (X^2 - c)(X^2 +
        c)$;
  \item $2 = c^2$: \quad $(X^2 + cX + 1)(X^2 - cX + 1) = X^4 + (2
        - c^2)X^2 + 1 = X^4 + 1$;
  \item $-2 = c^2$: \quad $(X^2 + cX - 1)(X^2 - cX - 1) = X^4 -
        (c^2 + 2)X^2 + 1 = X^4 + 1$.
\end{itemize}
So $X^4+1$ is reducible modulo every prime, yet irreducible over
$\Q$: the reduction criterion (\cref{thm:b3:rings:criteria}(1))
detects irreducibility but its failure proves nothing.

(For the structural reason: $p^2 - 1 = (p-1)(p+1)$ is a product of
two consecutive even numbers, so $8 \mid p^2 - 1$; the cyclic
group $\mathbb F_{p^2}^\times$ (cyclicity proved in
\cref{ch:b3:galois}) then contains an element $\zeta$ of order
$8$, a root of $X^4 + 1$; its minimal polynomial over $\mathbb
F_p$ divides $X^4+1$ and has degree $\leq 2$ --- $X^4+1$ can never
be irreducible mod $p$.)
\end{solution}

\begin{solution}{pb:b3:rings:1}
\textbf{1.} $N(z) = z\bar z$, so $N(zw) = zw\overline{zw} = z\bar
z\, w \bar w = N(z)N(w)$. If $uv = 1$: $N(u)N(v) = 1$ in $\N$, so
$N(u) = 1$, i.e.\ $u \in \{\pm 1, \pm \iu\}$; all four are units.
Brahmagupta: $(a^2+b^2)(c^2+d^2) = N\bigl((a + \iu b)(c + \iu
d)\bigr) = (ac - bd)^2 + (ad + bc)^2$.

\textbf{2.} If $z = ab$, then $N(z) = N(a)N(b)$ is prime, so
$N(a) = 1$ or $N(b) = 1$: one factor is a unit. As $N(z) > 1$,
$z$ is neither zero nor a unit: irreducible --- and prime, since
$\Z[\iu]$ is a UFD
(\cref{thm:b3:rings:euclideanpid,thm:b3:rings:pidufd,%
lem:b3:rings:bezout}).

\textbf{3.} $\pi$ divides $N(\pi) = \pi\bar\pi \geq 2$, an
integer; factoring $N(\pi)$ into prime numbers and using that
$\pi$ is prime, $\pi \mid p$ for some prime number $p$. If also
$\pi \mid q \neq p$: B\'ezout in $\Z$ gives $1 = up + vq$, so
$\pi \mid 1$ --- absurd: $p$ is unique. From $p = \pi\gamma$:
$p^2 = N(p) = N(\pi)N(\gamma)$ with $N(\pi) \neq 1$, so $N(\pi)
\in \{p, p^2\}$.

\textbf{4.} Let $\pi$ be a Gaussian prime dividing $p$, $p =
\pi\gamma$. If $N(\pi) = p^2$: $N(\gamma) = 1$, so $p$ is an
associate of $\pi$, itself a Gaussian prime; and no Gaussian
prime has norm $p$ (if $N(\rho) = p$ then $\rho \mid
\rho\bar\rho = p$, and $p$ prime in $\Z[\iu]$ would force
$\rho$ associate to $p$, giving $N(\rho) = N(p) = p^2 \neq
p$).
If $N(\pi) = p$: writing $\pi = a + \iu b$, $p = \pi\bar\pi = a^2
+ b^2$.

\textbf{5.} In the abelian group $(\Z/p\Z)^\times$, pair each
element with its inverse. The self-inverse elements are the roots
of $X^2 - 1$: exactly $\pm 1$ (at most two roots in a field).
The product of all elements is then $1 \cdot (-1) \cdot \prod
(\text{pairs } k k^{-1}) = -1$: $(p-1)! \equiv -1 \pmod p$. (For
$p = 2$: $1! \equiv -1$.)

\textbf{6.} Write $(p-1)! = \prod_{k=1}^m k \cdot
\prod_{k=m+1}^{p-1}k$ with $m = \frac{p-1}2$. In the second
product substitute $k = p - j$, $j = 1, \dots, m$: modulo $p$,
$\prod_{j=1}^m (p - j) \equiv (-1)^m m!$. Hence $-1 \equiv
(-1)^m (m!)^2$, i.e.\ $(m!)^2 \equiv (-1)^{m+1} \pmod p$.

\textbf{7.} If $p \equiv 1 \pmod 4$, $m$ is even and question 6
gives $(m!)^2 \equiv -1$: a square root of $-1$. Conversely, if
$x^2 \equiv -1 \pmod p$ ($p$ odd), then $x^4 = 1 \neq x^2$: $x$
has order $4$ in $(\Z/p\Z)^\times$, so $4 \mid p - 1$ (Lagrange).
And $p = 2$: $1^2 = 1 \equiv -1$. Conclusion: $-1$ is a square
mod $p$ iff $p = 2$ or $p \equiv 1 \pmod 4$.

\textbf{8.} With $x^2 \equiv -1$: $p \mid x^2 + 1 = (x + \iu)(x -
\iu)$. If $p$ were a Gaussian prime it would divide a factor; but
$\frac xp \pm \frac \iu p \notin \Z[\iu]$. So $p$ is not a
Gaussian prime; by the dichotomy (question 4) --- $p$ not prime
means the second branch --- $p = a^2 + b^2$.

\textbf{9.} Squares are $\equiv 0$ or $1 \pmod 4$, so $a^2 + b^2
\in \{0, 1, 2\} \pmod 4$: a prime $p \equiv 3 \pmod 4$ is not a
sum of two squares. By question 4, the branch $N(\pi) = p$ ($p =
a^2+b^2$) is impossible: $p$ stays a Gaussian prime.

\textbf{10.} $(1 + \iu)^2 = 2\iu$, so $2 = -\iu(1 + \iu)^2$; and
$N(1 + \iu) = 2$ is prime, so $1 + \iu$ is a Gaussian prime
(question 2).

\textbf{11.} Every Gaussian prime divides exactly one prime
number $p$ (question 3); listing by cases: $p = 2$ gives the
associates of $1 + \iu$; $p \equiv 3 \pmod 4$ gives $p$ itself
(question 9); $p \equiv 1 \pmod 4$ gives the pair $\pi, \bar\pi$
of norm $p$ (questions 4 and 8). The pair is genuine: $\bar\pi \in
\{\pm\pi, \pm\iu\pi\}$ would force, writing $\pi = a + \iu b$,
either $b = 0$, $a = 0$, or $a = \pm b$, giving $p = a^2 + b^2 \in
\{a^2, 2a^2\}$ --- impossible for an odd prime. Check:
$5 = (2 + \iu)(2 - \iu)$, $N(2\pm\iu) = 5$; $3$: prime of norm
$9$.

\textbf{12.} Write $n = 2^{\alpha}\prod_i p_i^{\beta_i} \prod_j
q_j^{2\gamma_j}$ with $p_i \equiv 1$, $q_j \equiv 3 \pmod 4$.
Each factor is a sum of two squares: $2 = 1^2 + 1^2$; $p_i = a^2
+ b^2$ (question 8); $q_j^{2\gamma_j} = (q_j^{\gamma_j})^2 +
0^2$. The Brahmagupta identity (question 1) propagates the
property to the product $n$.

\textbf{13.} Let $n = a^2 + b^2 = N(a + \iu b)$ and $q \equiv 3
\pmod 4$, $q \mid n$. The Gaussian prime $q$ (question 9) divides
$(a + \iu b)(a - \iu b)$, hence one of the two factors --- say $q
\mid a + \iu b$ (the other case is identical). But then $\frac{a +
\iu b}{q} = \frac aq + \iu \frac bq \in \Z[\iu]$ reads off as $q
\mid a$ \emph{and} $q \mid b$ in $\Z$. Hence $q^2 \mid n$ and
$\frac n{q^2} =
\bigl(\frac aq\bigr)^2 + \bigl(\frac bq\bigr)^2$. By strong
induction on $n$, the exponent of $q$ in $n/q^2$ is even; that of
$n$ is even too.

\textbf{14.} \emph{Theorem (Fermat).} A positive integer is a sum
of two squares if and only if every prime $\equiv 3 \pmod 4$
occurs in it with an even exponent. --- $2025 = 3^4 \cdot 5^2$:
exponent of $3$ even, yes ($2025 = 45^2 + 0^2 = 27^2 + 36^2$).
$2026 = 2 \cdot 1013$ with $1013 \equiv 1 \pmod 4$ prime: yes
($1013 = 22^2 + 23^2$, and Brahmagupta with $2 = 1^2+1^2$:
$2026 = (22 - 23)^2 + (22 + 23)^2 = 1^2 + 45^2$). $2027$ is a
prime $\equiv 3 \pmod 4$: no.

\textbf{15.} Let $p = a^2 + b^2 = c^2 + d^2$ with positive
integers, $p \equiv 1 \pmod 4$, and $\pi$ a Gaussian prime with
$p = \pi\bar\pi$ (question 4). Both $a + \iu b$ and $c + \iu d$
have norm $p$, hence are Gaussian primes (question 2) dividing
$p = (a+\iu b)(a - \iu b)$; by uniqueness of factorization, $c +
\iu d$ is an associate of $a + \iu b$ or of $a - \iu b$:
\[
c + \iu d \in \{\pm(a \pm \iu b),\ \pm\iu(a \pm \iu b)\}
= \{\pm a \pm \iu b,\ \pm b \pm \iu a\}.
\]
Positivity of $c, d$ leaves $c + \iu d \in \{a + \iu b, b + \iu
a\}$: $\{c, d\} = \{a, b\}$.
\end{solution}
